% Future Work. Moved out of the Conclusions chapter: the structure guide
% asks for a half-page Conclusions section and no Future Work chapter, so
% the material is kept here rather than discarded.

These directions follow from the thesis but were not pursued in it.

\textbf{Weighted aggregation up the guideline hierarchy.} Recommendations are currently averaged with equal weight, so a critical requirement counts the same as a peripheral one. The quality-assessment literature aggregates up a weighted tree mirroring the guideline's plan structure; adopting that would let the overall score reflect clinical priority rather than mere count.

\textbf{Cross-cohort pretraining.} Both Kineret models were trained from scratch so that architecture and initialization would not be confounded. A model pretrained across hospital networks and fine-tuned per site is the natural route to something robust rather than fitted to one source, and it would test whether the knowledge base's cohort-specific advantage survives transfer. Building a large pretrained artifact in the spirit of HEART \citep{huang2024heart} would open fine-tuning to other diseases and care settings, and would be a step toward a general-purpose clinical representation over interval data.

\textbf{Adherence as a model input at scale.} The null result here was measured on one cohort, at one evaluation window, with adherence entering as a mean per pattern. Richer encodings --- the per-day profile rather than its mean, or the compliance trajectory as a sequence --- were not tried, and whether a quality-of-care signal helps prediction is not settled by a single aggregation.

\textbf{Mining candidate recommendations.} Because each recommendation is declared externally as a pattern rather than implemented in analysis code, evaluating a practice merely \emph{suspected} to matter costs no more than evaluating the fourteen scored here. KarmaLego offers a route to finding such candidates rather than guessing them: mine frequent interval patterns, rank them by feature importance over the TPF representation, and translate the survivors into TAK definitions that can then be scored and tested like any other recommendation. That inverts the usual question --- which practices earn a place in a guideline, and which current recommendations track no outcome at all. Any candidate found this way is a hypothesis for prospective evaluation, and a wide search would need more conservative multiplicity correction than was applied here.

\textbf{Interval representation learning beyond healthcare.} INTERVenE is built around properties of interval sequences --- that a symbol holds over a span, that spans overlap --- rather than around anything specific to medicine. Comparing interval representations systematically, and testing the architecture on another interval domain such as industrial failure prediction or process monitoring, would move the claim from architecture engineering toward the representation question: why interval representations work, and when they are worth their cost.

\textbf{Explainability as the primary object.} The attribution work here is a demonstration rather than a study. A systematic comparison of explanation quality between a semantic interval vocabulary and a non-semantic one --- the same $\sigma$-bin contrast, judged on interpretability rather than accuracy --- would test the central claim directly: that semantic vocabularies produce clinically meaningful explanations at no cost in predictive performance.

\textbf{Prospective evaluation.} Everything in this thesis is retrospective. The measurement instrument is automated, component-level, and transferable, but whether surfacing these scores to a treating team changes practice --- and whether changed practice changes outcomes --- is a question only a prospective study can answer.
